\documentclass[12pt,a4paper]{article}

\usepackage[utf8]{inputenc}
\usepackage[T1]{fontenc}
\usepackage[spanish]{babel}
\usepackage{amsmath,amssymb}
\usepackage{array,tabularx}
\usepackage{tikz}
\usetikzlibrary{arrows.meta}
\usepackage[margin=2cm]{geometry}

\title{Práctica: caracterización de detectores CCD}
\author{\mbox{}}
\date{}

\begin{document}

\maketitle

\section{Caracterización de la ganancia y el ruido de lectura}

La ganancia, $G$, es el factor de conversión entre la carga generada en el
semiconductor y las unidades digitales (ADU) registradas por el convertidor
análogo-digital. Se expresa en electrones por ADU, $\mathrm{e^-}/\mathrm{ADU}$.
El ruido de lectura, $\epsilon_{\text{RN}}$, cuantifica la incertidumbre que
introduce la electrónica al medir la carga de cada píxel.

\begin{figure}[htbp]
    \centering
    \begin{tikzpicture}[
        font=\small,
        luz/.style={-{Stealth[length=2mm]}, draw=orange!80!black,
            line width=0.9pt},
        estructura/.style={draw=gray!75!black, line width=1.1pt},
        etiqueta/.style={font=\scriptsize\bfseries, align=center,
            text=blue!60!black},
        datos/.style={-{Stealth[length=2mm]}, draw=teal!70!black,
            dashed, line width=1pt}
    ]
        \node[font=\bfseries\large, text=blue!65!black] at (8.1,4.2)
            {Montaje experimental aproximado};

        % Trípode y montura del telescopio
        \draw[estructura] (7.45,1.82) -- (7.45,1.1);
        \draw[estructura] (7.45,1.1) -- (6.45,0.2);
        \draw[estructura] (7.45,1.1) -- (8.45,0.2);
        \draw[estructura] (7.45,1.1) -- (7.45,0.2);
        \draw[estructura, line width=2pt] (6.15,0.18) -- (6.7,0.18);
        \draw[estructura, line width=2pt] (8.2,0.18) -- (8.75,0.18);
        \draw[estructura, line width=2pt] (7.2,0.18) -- (7.7,0.18);
        \draw[fill=gray!25, draw=gray!70!black] (7.05,1.52)
            rectangle (7.85,1.95);
        \draw[fill=gray!45, draw=gray!70!black] (7.28,1.83)
            circle (0.28);

        % Pantalla utilizada como fuente de iluminación
        \draw[fill=gray!25, draw=gray!65!black, line width=1pt,
            rounded corners=2pt] (0.45,1.55) rectangle (2.95,3.3);
        \draw[fill=yellow!5, draw=blue!45!black, line width=0.7pt]
            (0.6,1.7) rectangle (2.8,3.15);
        \node[font=\scriptsize\bfseries, text=gray!65!black]
            at (1.7,2.55) {BLANCO};
        \node[font=\tiny, text=gray!55!black] at (1.7,2.3)
            {campo uniforme};
        \draw[fill=gray!55, draw=gray!70!black]
            (1.52,1.55) -- (1.88,1.55) -- (2.02,0.92) --
            (1.38,0.92) -- cycle;
        \draw[fill=gray!35, draw=gray!70!black, rounded corners=1pt]
            (1.08,0.78) rectangle (2.32,0.92);

        % Luz de la pantalla que entra por el objetivo
        \fill[orange!18, opacity=0.75]
            (2.82,1.78) -- (4.35,1.9) -- (4.35,2.92) --
            (2.82,3.03) -- cycle;
        \draw[luz] (2.82,2.95) -- (4.28,2.78);
        \draw[luz] (2.82,2.38) -- (4.28,2.38);
        \draw[luz] (2.82,1.83) -- (4.28,2.0);

        % Tubo, anillos y objetivo
        \draw[fill=blue!9, draw=blue!55!black, line width=1.2pt,
            rounded corners=7pt] (4.35,1.83) rectangle (10.25,2.95);
        \draw[fill=blue!18, draw=blue!55!black]
            (4.8,1.88) rectangle (9.83,2.9);
        \draw[fill=gray!35, draw=gray!70!black]
            (5.05,1.82) rectangle (5.25,2.96);
        \draw[fill=gray!35, draw=gray!70!black]
            (9.72,1.82) rectangle (9.92,2.96);
        \draw[fill=gray!30, draw=gray!65!black, line width=1pt]
            (4.25,1.72) rectangle (4.58,3.06);
        \draw[fill=cyan!25, draw=blue!65!black, opacity=0.8]
            (4.26,1.83) ellipse [x radius=0.08, y radius=0.55];

        % Convergencia de los rayos en el plano del detector
        \draw[luz] (4.42,2.78) -- (10.55,2.4);
        \draw[luz] (4.42,2.38) -- (10.55,2.4);
        \draw[luz] (4.42,2.0) -- (10.55,2.4);

        % Enfocador y cámara CCD acoplada al extremo posterior
        \draw[fill=gray!35, draw=gray!70!black, line width=1pt]
            (10.15,2.0) rectangle (10.65,2.78);
        \draw[fill=gray!55, draw=gray!70!black]
            (10.55,2.12) rectangle (10.78,2.66);
        \draw[fill=violet!22, draw=violet!65!black, line width=1pt,
            rounded corners=2pt] (10.78,1.85) rectangle (12.05,2.94);
        \draw[fill=violet!38, draw=violet!65!black]
            (11.0,1.95) rectangle (11.82,2.12);
        \draw[fill=cyan!45, draw=teal!65!black, line width=0.7pt]
            (10.8,2.18) rectangle (10.92,2.6);
        \draw[draw=violet!65!black, line width=0.8pt]
            (11.1,2.7) -- (11.72,2.7);

        % Computador que controla la cámara y recibe los datos
        \draw[fill=gray!30, draw=gray!70!black, line width=1pt,
            rounded corners=2pt] (13.35,1.08) rectangle (15.35,2.48);
        \draw[fill=teal!8, draw=teal!65!black]
            (13.48,1.22) rectangle (15.22,2.34);
        \node[font=\tiny\bfseries, text=teal!60!black, align=center]
            at (14.35,1.82) {CONTROL Y\\LECTURA CCD};
        \draw[fill=gray!40, draw=gray!70!black]
            (14.18,1.08) rectangle (14.52,0.82);
        \draw[fill=gray!28, draw=gray!70!black, rounded corners=1pt]
            (13.82,0.68) rectangle (14.88,0.82);
        \draw[datos] (11.55,1.88) .. controls (12.25,0.8) and
            (12.7,0.78) .. (13.35,1.35);
        \node[font=\tiny, text=teal!60!black] at (12.55,0.62) {USB};

        % Etiquetas con líneas guía
        \node[etiqueta] at (1.7,3.62) {Computador como\\fuente flat};
        \draw[-{Stealth[length=1.5mm]}, draw=blue!55!black]
            (1.7,3.37) -- (1.7,3.16);
        \node[etiqueta] at (7.45,3.53) {Telescopio};
        \draw[-{Stealth[length=1.5mm]}, draw=blue!55!black]
            (7.45,3.35) -- (7.45,2.98);
        \node[etiqueta] at (11.42,3.53) {Cámara CCD};
        \draw[-{Stealth[length=1.5mm]}, draw=blue!55!black]
            (11.42,3.35) -- (11.42,2.96);
        \node[etiqueta] at (14.35,2.78) {Computador de adquisición};
        \draw[-{Stealth[length=1.5mm]}, draw=blue!55!black]
            (14.35,2.63) -- (14.35,2.49);

        \draw[gray!65, line width=0.7pt] (0.15,0.16) -- (15.8,0.16);
    \end{tikzpicture}
    \caption{Vista lateral ilustrativa del montaje: una pantalla que muestra
    un campo blanco uniforme ilumina el objetivo del telescopio; la cámara
    CCD se acopla al extremo posterior y se conecta al computador de
    adquisición. Para tomar bias, bloquee la luz antes del detector.}
    \label{fig:montaje-ccd}
\end{figure}
\begin{figure}[htbp]
    \centering
    \begin{tikzpicture}[
        font=\small,
        region/.style={draw=red!75!black, dashed, line width=1pt},
        borde/.style={draw=gray!65!black, line width=0.7pt}
    ]
        \node[font=\bfseries\large, text=blue!65!black] at (8.5,4.3)
            {Ejemplos de imágenes para estimar $G$ y $\epsilon_{\text{RN}}$};

        \node[font=\small\bfseries] at (2.75,3.78) {$\mathrm{Flat}_1$};
        \node[font=\small\bfseries] at (8.5,3.78) {$\mathrm{Flat}_2$};
        \node[font=\small\bfseries] at (14.25,3.78)
            {$\Delta_F=\mathrm{Flat}_1-\mathrm{Flat}_2$};
        \node[font=\small\bfseries] at (2.75,1.63) {$\mathrm{Bias}_1$};
        \node[font=\small\bfseries] at (8.5,1.63) {$\mathrm{Bias}_2$};
        \node[font=\small\bfseries] at (14.25,1.63)
            {$\Delta_B=\mathrm{Bias}_1-\mathrm{Bias}_2$};

        % Matrices esquemáticas: tonos uniformes para frames originales
        \foreach \xa/\xb/\ya/\yb/\tono in {
            0.4/5.1/2.0/3.5/gray!12,
            6.15/10.85/2.0/3.5/gray!14,
            11.9/16.6/2.0/3.5/gray!65,
            0.4/5.1/-0.15/1.35/gray!42,
            6.15/10.85/-0.15/1.35/gray!45,
            11.9/16.6/-0.15/1.35/gray!65
        } {
            \fill[\tono] (\xa,\ya) rectangle (\xb,\yb);
            \draw[step=0.4cm, draw=gray!35, line width=0.2pt]
                (\xa,\ya) grid (\xb,\yb);
            \draw[borde] (\xa,\ya) rectangle (\xb,\yb);
        }

        % Fluctuaciones de píxel a píxel en las imágenes diferencia
        \foreach \x/\y in {
            12.2/2.25, 12.85/2.95, 13.45/2.5, 14.15/3.12,
            14.85/2.35, 15.55/2.78, 16.05/3.08
        } {
            \fill[gray!16] (\x,\y) rectangle ++(0.25,0.22);
        }
        \foreach \x/\y in {
            12.65/2.62, 13.2/2.18, 13.85/2.83, 14.55/2.62,
            15.25/3.02, 15.85/2.3
        } {
            \fill[gray!92] (\x,\y) rectangle ++(0.25,0.22);
        }
        \foreach \x/\y in {12.55/0.28, 13.6/0.88, 14.85/0.42, 15.75/0.92} {
            \fill[gray!35] (\x,\y) rectangle ++(0.25,0.22);
        }
        \foreach \x/\y in {13.05/0.72, 14.1/0.25, 15.35/0.8} {
            \fill[gray!88] (\x,\y) rectangle ++(0.25,0.22);
        }

        % Misma región central de medida en cada frame
        \foreach \x in {2.0,7.75,13.5} {
            \draw[region] (\x,2.38) rectangle ++(1.5,0.72);
            \draw[region] (\x,0.23) rectangle ++(1.5,0.72);
        }
    \end{tikzpicture}
    \caption{Representación ilustrativa, no basada en datos reales, de dos
    flats y dos bias junto con sus imágenes diferencia. El contorno discontinuo
    señala la misma región central usada para medir medias y dispersiones;
    las diferencias muestran las fluctuaciones píxel a píxel.}
    \label{fig:imagenes-flat-bias}
\end{figure}

\subsection{Adquisición y preparación de las imágenes}

\begin{enumerate}
    \item Tome una secuencia de 40 imágenes bias con la cámara CCD.
    \item Tome una serie de imágenes \emph{flat-field} con distintos tiempos
    de exposición. Comience en $0.05\,\mathrm{s}$ y duplique el tiempo en cada
    paso ($0.05$, $0.1$, $0.2$, $0.4\,\mathrm{s}$, etc.). Seleccione imágenes
    bien iluminadas (alrededor de 15\,000 ADU o más), pero evite las zonas
    saturadas o no lineales.
    \item Para cada tiempo de exposición, divida los flats en dos grupos y
    combine cada grupo por separado con \texttt{imcombine} y
    \texttt{combine=average}. Las imágenes resultantes se denotarán
    $\text{Flat}_1$ y $\text{Flat}_2$. Tome también dos imágenes bias,
    $\text{Bias}_1$ y $\text{Bias}_2$.
    \item Reste las imágenes de cada par para obtener
    \begin{align*}
        \Delta_F &= \text{Flat}_1 - \text{Flat}_2,\\
        \Delta_B &= \text{Bias}_1 - \text{Bias}_2.
    \end{align*}
    Calcule las medias de cada imagen y las desviaciones estándar de
    $\Delta_F$ y $\Delta_B$ en una región central, evitando los bordes.
\end{enumerate}

Por ejemplo, puede dividir en dos grupos los flats tomados con el mismo tiempo
de exposición (como $2\,\mathrm{s}$) y combinar cada grupo por promedio. Use
ambas imágenes resultantes para calcular las medias y su imagen diferencia.

\subsection{Cálculo de la ganancia y el ruido de lectura}

Calcule la ganancia para cada tiempo de exposición mediante
\begin{equation}
    G = \frac{(\bar{F}_1 + \bar{F}_2) - (\bar{B}_1 + \bar{B}_2)}
    {\sigma^2_{\Delta_F} - \sigma^2_{\Delta_B}},
    \label{eq:gain}
\end{equation}
donde $\bar{F}_1$ y $\bar{F}_2$ son las medias de los dos flats,
$\bar{B}_1$ y $\bar{B}_2$ son las medias de los dos bias, y
$\sigma^2_{\Delta_F}$ y $\sigma^2_{\Delta_B}$ son las varianzas de las
imágenes diferencia respectivas.

Una vez conocida la ganancia, estime el ruido de lectura con
\begin{equation}
    \epsilon_{\text{RN}} = G\,\frac{\sigma_{\Delta_B}}{\sqrt{2}}.
    \label{eq:read-noise}
\end{equation}

Registre los datos para cada tiempo de exposición:

\begin{center}
\scriptsize
\setlength{\tabcolsep}{3pt}
\renewcommand{\arraystretch}{1.3}
\begin{tabularx}{\textwidth}{|c|*{7}{>{\centering\arraybackslash}X|}}
    \hline
    \textbf{Tiempo} &
    \shortstack{\textbf{Media}\\\textbf{Flat 1}\\\textbf{(ADU)}} &
    \shortstack{\textbf{Media}\\\textbf{Flat 2}\\\textbf{(ADU)}} &
    \shortstack{\textbf{Media}\\\textbf{Bias 1}\\\textbf{(ADU)}} &
    \shortstack{\textbf{Media}\\\textbf{Bias 2}\\\textbf{(ADU)}} &
    \shortstack{$\boldsymbol{\sigma_{\Delta_F}}$\\\textbf{(ADU)}} &
    \shortstack{$\boldsymbol{\sigma_{\Delta_B}}$\\\textbf{(ADU)}} &
    \shortstack{\textbf{$G$}\\\textbf{(e$^-$/ADU)}} \\
    \hline
    $t_{\text{exp}}$ (s) & & & & & & & \\
    \hline
    0.05 & & & & & & & \\
    \hline
    0.10 & & & & & & & \\
    \hline
    $\vdots$ & & & & & & & \\
    \hline
\end{tabularx}
\end{center}

\subsection{Análisis y discusión}

\begin{itemize}
    \item Grafique $G$ en función del tiempo de exposición de los flats.
    ¿Qué comportamiento observa?
    \item Grafique $\epsilon_{\text{RN}}$ en función del tiempo de exposición
    usado para calcular $G$ y presente una tabla con los resultados.
    \item Investigue qué hace la tarea \texttt{findgain} de IRAF. Ejecútela y
    compare su resultado con los obtenidos en esta práctica.
    \item Busque los valores de $G$ y $\epsilon_{\text{RN}}$ reportados por el
    fabricante de la cámara. Compare los resultados y discuta las posibles
    diferencias.
\end{itemize}

\section{Respuesta lineal del detector}

Un detector ideal responde de forma lineal: las cuentas registradas (ADU) son
directamente proporcionales al número de fotones incidentes y al tiempo de
exposición. A niveles altos de iluminación, los pozos de potencial se llenan
hasta alcanzar la saturación y la respuesta deja de ser lineal.

\begin{figure}[htbp]
    \centering
    \begin{tikzpicture}[
        font=\scriptsize,
        ideal/.style={draw=gray!75!black, dashed, line width=1.2pt},
        medida/.style={draw=blue!75!black, line width=1.8pt}
    ]
        \node[font=\bfseries\large, text=blue!65!black] at (5.8,5.35)
            {Respuesta media del detector frente al tiempo de exposición};

        \fill[green!8] (0.8,0.8) rectangle (5.8,4.55);
        \fill[yellow!14] (5.8,0.8) rectangle (8.0,4.55);
        \fill[red!7] (8.0,0.8) rectangle (10.8,4.55);
        \draw[step=1cm, draw=gray!28, line width=0.25pt]
            (0.8,0.8) grid (10.8,4.55);

        \draw[-{Stealth[length=2mm]}, line width=0.9pt]
            (0.8,0.8) -- (11.0,0.8);
        \draw[-{Stealth[length=2mm]}, line width=0.9pt]
            (0.8,0.8) -- (0.8,4.75);

        \foreach \x/\t in {0/0,2/2,4/4,6/6,8/8,10/10} {
            \draw (0.8+\x,0.8) -- (0.8+\x,0.72);
            \node[below] at (0.8+\x,0.68) {$\t$};
        }
        \foreach \y/\s in {0/0,1/1,2/2,3/3,4/4,5/5} {
            \draw (0.8,0.8+0.75*\y) -- (0.72,0.8+0.75*\y);
            \node[left] at (0.68,0.8+0.75*\y) {$\s$};
        }

        \node[align=center, text=green!40!black] at (3.25,4.72)
            {Régimen\\lineal};
        \node[align=center, text=orange!60!black] at (6.9,4.72)
            {Pérdida de\\linealidad};
        \node[align=center, text=red!65!black] at (9.4,4.72)
            {Saturación};

        \draw[ideal] (0.8,0.8) -- (10.5,4.44);
        \draw[medida] plot[smooth, tension=0.8] coordinates {
            (0.8,0.8) (2.8,1.55) (4.8,2.3) (5.8,2.68)
            (6.8,3.12) (7.8,3.44) (8.8,3.52) (10.5,3.53)
        };
        \node[fill=white, inner sep=1.5pt, text=blue!70!black]
            at (3.9,2.55) {respuesta medida};
        \node[fill=white, inner sep=1.5pt, text=gray!70!black]
            at (8.15,4.08) {proporcionalidad ideal};

        \node[align=center] at (5.8,0.12)
            {Tiempo de exposición (s)};
        \node[rotate=90, align=center] at (0.05,2.7)
            {Señal media (ADU)};
    \end{tikzpicture}
    \caption{Curva conceptual: la señal crece proporcionalmente durante el
    régimen lineal, se desvía de esa tendencia y finalmente alcanza la
    saturación. Reemplace la curva ilustrativa por los datos medidos.}
    \label{fig:linealidad-detector}
\end{figure}
\begin{figure}[htbp]
    \centering
    \begin{tikzpicture}[
        font=\small,
        borde/.style={draw=gray!65!black, line width=0.45pt},
        rebose/.style={-{Stealth[length=1.8mm]}, draw=blue!70!black,
            dashed, line width=1pt}
    ]
        \node[font=\bfseries\large, text=blue!65!black] at (8.1,5.25)
            {Carga almacenada en los píxeles de un CCD};
        \node[font=\bfseries, text=green!40!black] at (2.9,4.72)
            {Régimen no saturado};
        \node[font=\bfseries, text=red!65!black] at (13.1,4.72)
            {Saturación y blooming};

        % Cada matriz tiene cinco filas y cinco columnas de igual tamaño.
        \fill[gray!6] (1.1,0.75) rectangle (4.7,4.35);
        \fill[yellow!35] (2.54,2.19) rectangle (3.26,2.91);
        \draw[step=0.72cm, borde] (1.1,0.75) grid (4.7,4.35);
        \draw[orange!75!black, line width=1pt]
            (2.54,2.19) rectangle (3.26,2.91);
        \foreach \x/\y in {
            2.72/2.38, 3.02/2.38, 2.72/2.67, 3.02/2.67
        } {
            \fill[orange!90!black] (\x,\y) circle (0.055);
        }
        \node[font=\scriptsize, align=center] at (2.9,0.38)
            {$Q<Q_{\mathrm{sat}}$: carga confinada al píxel};

        \fill[gray!6] (11.3,0.75) rectangle (14.9,4.35);
        \fill[orange!15] (12.74,3.63) rectangle (13.46,4.35);
        \fill[orange!45] (12.74,2.91) rectangle (13.46,3.63);
        \fill[red!75!black] (12.74,2.19) rectangle (13.46,2.91);
        \fill[orange!45] (12.74,1.47) rectangle (13.46,2.19);
        \fill[orange!15] (12.74,0.75) rectangle (13.46,1.47);
        \draw[step=0.72cm, borde] (11.3,0.75) grid (14.9,4.35);
        \draw[red!60!black, line width=1.1pt]
            (12.74,2.19) rectangle (13.46,2.91);
        \foreach \x/\y in {
            12.88/2.34, 13.12/2.34, 12.88/2.55, 13.12/2.55,
            12.88/2.75, 13.12/2.75
        } {
            \fill[white] (\x,\y) circle (0.05);
        }
        \draw[rebose] (13.1,2.91) -- (13.1,3.5);
        \draw[rebose] (13.1,2.19) -- (13.1,1.6);
        \node[font=\scriptsize, align=center] at (13.1,0.38)
            {$Q\geq Q_{\mathrm{sat}}$: el exceso alcanza píxeles vecinos};

        % Leyenda: los puntos representan electrones acumulados.
        \draw[fill=gray!6, borde] (0.75,-0.18) rectangle (1.05,0.12);
        \node[font=\tiny, anchor=west] at (1.13,-0.03) {píxel sin carga};
        \draw[fill=yellow!35, borde] (4.25,-0.18) rectangle (4.55,0.12);
        \fill[orange!90!black] (4.4,-0.03) circle (0.045);
        \node[font=\tiny, anchor=west] at (4.63,-0.03)
            {electrones almacenados};
        \draw[fill=red!75!black, borde] (8.35,-0.18) rectangle (8.65,0.12);
        \fill[white] (8.5,-0.03) circle (0.045);
        \node[font=\tiny, anchor=west] at (8.73,-0.03) {pozo saturado};
        \draw[fill=orange!45, borde] (11.8,-0.18) rectangle (12.1,0.12);
        \node[font=\tiny, anchor=west] at (12.18,-0.03)
            {carga desbordada};
    \end{tikzpicture}
    \caption{Ilustración conceptual del blooming. En el régimen no saturado,
    la carga se mantiene dentro del píxel iluminado. Al alcanzar su capacidad,
    el exceso puede pasar a los píxeles vecinos y producir una traza a lo largo
    de una columna; el patrón exacto depende del detector.}
    \label{fig:blooming-detector}
\end{figure}

\begin{itemize}
    \item Incremente progresivamente el tiempo de exposición de los flats
    hasta observar la saturación. En esa región, la señal deja de crecer
    linealmente, se estabiliza o alcanza valores límite del formateador.
    \item Grafique el flujo promedio de los píxeles de cada flat-field en
    función del tiempo de exposición.
    \item Determine numéricamente el límite del régimen lineal de la cámara,
    en ADU, y calcule el rango dinámico del detector.
    \item Discuta el fenómeno de \textit{blooming} (desborde de carga), que
    ocurre cuando un píxel supera su capacidad de almacenamiento.
\end{itemize}

\section{Tasa de corriente oscura}

La corriente oscura es la tasa de acumulación de electrones generados por
agitación térmica en el silicio, aun en ausencia de luz. Se expresa en
$\mathrm{e^-\,pixel^{-1}\,s^{-1}}$ o
$\mathrm{ADU\,pixel^{-1}\,s^{-1}}$ para una temperatura de operación dada.
\begin{figure}[htbp]
    \centering
    \begin{tikzpicture}[
        font=\small,
        carcasa/.style={draw=gray!75!black, line width=1pt},
        etiqueta/.style={font=\scriptsize\bfseries, text=blue!60!black,
            align=center},
        conexion/.style={draw=teal!70!black, line width=1.1pt}
    ]
        \node[font=\bfseries\large, text=blue!65!black] at (8,5.35)
            {Adquisición de imágenes dark};

        % Cámara CCD con la entrada óptica bloqueada por una tapa opaca.
        \draw[fill=gray!18, carcasa, rounded corners=5pt]
            (2.1,2.05) rectangle (7.45,4.05);
        \draw[fill=gray!35, carcasa] (2.45,2.3) rectangle (2.85,3.8);
        \foreach \y in {2.48,2.76,3.04,3.32,3.6} {
            \draw[draw=gray!65] (3.15,\y) -- (3.5,\y);
        }
        \draw[fill=gray!45, carcasa] (7.45,2.35) rectangle (7.78,3.75);
        \draw[fill=violet!22, draw=violet!65!black, line width=0.9pt]
            (7.78,2.25) rectangle (8.45,3.85);
        \draw[fill=cyan!45, draw=teal!65!black, line width=0.8pt]
            (7.88,2.55) rectangle (8.02,3.55);
        \node[font=\tiny\bfseries, text=violet!65!black, align=center]
            at (8.25,3.04) {CCD};
        \draw[fill=gray!40, carcasa] (3.25,1.72) rectangle (6.25,2.05);
        \node[etiqueta] at (5,4.4) {Cámara CCD};

        % Tapa frontal: la cara oscura representa una superficie opaca.
        \draw[fill=gray!15, carcasa, line width=1.2pt]
            (1.45,2.55) rectangle (2.1,3.55);
        \draw[fill=black!85, draw=black, line width=0.8pt]
            (1.25,2.48) ellipse (0.18 and 0.56);
        \node[etiqueta] at (1.8,4.4) {Tapa opaca\\o shutter cerrado};
        \draw[-{Stealth[length=1.7mm]}, draw=blue!55!black]
            (1.8,4.1) -- (1.8,3.62);

        % Computadora de control con una secuencia repetida de exposiciones.
        \draw[fill=gray!28, carcasa, rounded corners=2pt]
            (10.25,1.95) rectangle (14.75,4.2);
        \draw[fill=blue!4, draw=gray!65!black, line width=0.8pt]
            (10.48,2.2) rectangle (14.52,3.95);
        \node[font=\scriptsize\bfseries, text=blue!65!black]
            at (12.5,3.68) {SERIE DE EXPOSICIONES};
        \foreach \ybottom/\ytop/\tiempo in {
            3.10/3.40/10, 2.67/2.97/30, 2.24/2.54/60
        } {
            \draw[fill=teal!12, draw=teal!50!black, rounded corners=1pt]
                (11.05,\ybottom) rectangle (13.95,\ytop);
            \node[font=\scriptsize, anchor=west] at (11.25,\ybottom+0.15)
                {$t_{\mathrm{exp}}=\tiempo\,\mathrm{s}$};
            \node[font=\scriptsize\bfseries, text=teal!65!black,
                anchor=east] at (13.75,\ybottom+0.15) {$\times 3$};
        }
        \draw[fill=gray!50, carcasa] (12.15,1.68) rectangle (12.85,1.95);
        \draw[fill=gray!35, carcasa, rounded corners=1pt]
            (11.55,1.55) rectangle (13.45,1.68);
        \node[etiqueta] at (12.5,4.52) {Computador de adquisición};

        % Cable de datos; no representa el trayecto de la luz.
        \draw[conexion] (8.45,2.35) .. controls (9.1,1.25) and
            (9.55,1.25) .. (10.25,2.2);
        \node[font=\tiny, text=teal!65!black] at (9.35,1.25) {USB / datos};

        \node[draw=orange!65!black, fill=orange!7, rounded corners=2pt,
            text width=13.8cm, align=center, inner sep=5pt,
            font=\scriptsize] at (8,0.55)
            {Mantenga bloqueada toda la luz y constante la temperatura.
            Repita cada tiempo de exposición al menos tres veces para poder
            combinar los darks y reducir el ruido aleatorio.};
    \end{tikzpicture}
    \caption{Esquema ilustrativo de la toma de darks: la entrada óptica se
    mantiene cubierta o con el shutter cerrado mientras la cámara integra
    durante distintos tiempos. La secuencia mostrada es solo un ejemplo.}
    \label{fig:adquisicion-darks}
\end{figure}

\begin{enumerate}
    \item Tome al menos tres imágenes dark para cada tiempo de exposición,
    siguiendo el esquema usado para los flats y extendiendo las exposiciones
    hasta unos 2 o 5 minutos.
    \item Reste la señal bias de cada imagen dark.
    \item Combine con \texttt{imcombine} y \texttt{combine=average} los darks
    del mismo tiempo de exposición para construir imágenes \emph{superdark} y
    reducir el ruido aleatorio de lectura.
    \item Calcule la señal media de cada \emph{superdark} y grafíquela en
    función del tiempo de exposición, en ADU o electrones.
    \item Ajuste una recta de regresión lineal a los datos. ¿Qué significado
    físico tiene la pendiente?
    \item Convierta la pendiente a $\mathrm{e^-\,pixel^{-1}\,s^{-1}}$ usando
    la ganancia obtenida previamente. Discuta la importancia del enfriamiento
    (entre $-70\,^{\circ}\mathrm{C}$ y $-140\,^{\circ}\mathrm{C}$) en los CCD
    astronómicos.
\end{enumerate}

\begin{figure}[htbp]
    \centering
    \begin{tikzpicture}[
        font=\small,
        rejilla/.style={draw=gray!23, line width=0.35pt},
        ajuste/.style={draw=blue!75!black, line width=1.5pt},
        dato/.style={draw=orange!75!black, fill=orange!80!black}
    ]
        \node[font=\bfseries\large, text=blue!65!black] at (7.9,5.15)
            {La pendiente mide la acumulación de señal oscura};

        % Cuadrícula y ejes del gráfico conceptual.
        \draw[step=1cm, rejilla] (1.35,1.05) grid (10.55,4.35);
        \draw[-{Stealth[length=2mm]}, line width=0.9pt]
            (1.35,1.05) -- (10.8,1.05);
        \draw[-{Stealth[length=2mm]}, line width=0.9pt]
            (1.35,1.05) -- (1.35,4.62);

        \node[align=center, font=\scriptsize] at (5.95,0.48)
            {Tiempo de exposición, $t_{\mathrm{exp}}$ (s)};
        \node[rotate=90, align=center, font=\scriptsize] at (0.42,2.7)
            {Media del superdark, $\bar{D}$ (ADU/píxel)};

        % Puntos ilustrativos y ajuste lineal conceptual tras restar bias.
        \draw[ajuste] (1.52,1.12) -- (10.25,4.18);
        \foreach \x/\y in {
            2.1/1.34, 3.55/1.79, 5.05/2.31, 6.48/2.76,
            8.05/3.34, 9.55/3.82
        } {
            \filldraw[dato] (\x,\y) circle (2.4pt);
        }
        \node[fill=white, inner sep=2pt, font=\scriptsize,
            text=blue!70!black] at (7.95,2.2) {ajuste lineal};

        % Triángulo de pendiente asociado a dos puntos sobre el ajuste.
        \draw[draw=orange!80!black, dashed, line width=0.9pt]
            (3.55,1.83) -- (8.05,1.83) -- (8.05,3.41);
        \node[font=\scriptsize, text=orange!65!black]
            at (5.8,1.57) {$\Delta t$};
        \node[font=\scriptsize, text=orange!65!black, rotate=90]
            at (8.31,2.62) {$\Delta\bar{D}$};

        \node[draw=teal!60!black, fill=teal!5, rounded corners=2pt,
            text width=4.25cm, align=left, inner sep=6pt,
            font=\scriptsize] at (13.25,3.45)
            {Tras restar el bias, la pendiente del ajuste es la tasa media
            de señal oscura en ADU por píxel y segundo.};
        \node[draw=orange!65!black, fill=orange!7, rounded corners=2pt,
            text width=4.25cm, align=center, inner sep=6pt,
            font=\small] at (13.25,1.55)
            {$m=\dfrac{\Delta\bar{D}}{\Delta t}$\\[3pt]
            $i_{\mathrm{dark}}=G\,m$};
        \node[font=\tiny, align=center, text=gray!65!black] at (8,0.03)
            {Puntos y valores conceptuales; no corresponden a mediciones reales.};
    \end{tikzpicture}
    \caption{Ejemplo conceptual del ajuste de la señal media de los
    superdarks, una vez restada la señal bias, en función del tiempo de
    exposición. La pendiente $m$ se convierte a corriente oscura en
    $\mathrm{e^-\,pixel^{-1}\,s^{-1}}$ multiplicándola por la ganancia $G$.}
    \label{fig:pendiente-corriente-oscura}
\end{figure}
\section{Detección de defectos cosméticos}

Las matrices CCD reales presentan imperfecciones de fabricación o degradación
por radiación, como píxeles muertos (sensibilidad nula), píxeles calientes
(alta corriente oscura) y columnas defectuosas o trampas de carga.

\begin{itemize}
    \item ¿Cómo puede usar las imágenes \emph{flat-field} y \emph{superdark}
    recolectadas para encontrar automáticamente los píxeles muertos y
    calientes?
    \item ¿Cómo identificar otros defectos cosméticos, como columnas malas o
    trampas de carga? Presente capturas de pantalla de DS9 en las que marque
    estos defectos.
    \item Investigue para qué sirve la tarea \texttt{ccdmask} de IRAF. Describa
    cómo genera un mapa de máscara de píxeles defectuosos que pueda usarse al
    corregir imágenes científicas.
\end{itemize}

\end{document}